\chapter{问题描述与需求分析}
\section{问题背景与数据约束}
排序算法必须在内存中比较关键字，但实际文件可能远大于可用于排序的工作区。若把全部输入装入\code{vector}再调用内排序，空间需求随$N$增长，无法说明外存管理机制。本实验把原始数据保存在文件中，排序过程中每次只处理限定数量的整数，再把阶段结果写回文件。

外存算法的代价不仅是比较次数。逐个整数调用文件接口会产生大量细碎访问；将相邻值按块传输可以利用顺序访问，把一次I/O服务于多次比较。缓冲区因此承担两个不同粒度之间的衔接：文件按块读写，计算按元素取值或输出。操作系统仍可能缓存这些调用，故程序计数不等于硬盘实际寻道或物理读写次数。

归并段是文件中一个连续、内部非降序的整数序列。多个段分别有序，并不意味着按段号拼接后整体有序。算法先生成段，再反复合并，直到只剩一个逻辑最终段。实现保留旧段以便核验，因此“只剩一个”指当前参与归并的段号列表只剩一个，而非磁盘目录只剩一项。

\begin{table}[htbp]\centering\caption{全篇符号与单位}\label{tab:symbols}
\begin{tabularx}{\linewidth}{lY}\toprule
符号 & 含义\\\midrule
$N$ & 原始输入的整数总数，每个元素为\code{int32_t}。\\
$M$ & 分块内排或置换选择候选容量，以元素个数计。\\
$B$ & 一个数据块可容纳的整数数，常见平台字节容量为$4B$。\\
$k$ & 每组K路归并允许的最多路数，至少为2。\\
$R$ & 生成阶段产生的初始段数；$L_i$表示第$i$段长度。\\
$C$ & 二叉计划中所有内部节点长度之和，即理论合并输出元素量。\\
$S$ & 分阶段并行任务的同步阶段数量，不等同于归并轮数。\\\bottomrule
\end{tabularx}\end{table}

如表\ref{tab:symbols}所示，$M$不是整个进程的字节内存上限。缓冲区、败者树状态、目录、计划以及测试参考数组另占空间。正式排序入口\code{sortFile}接收文件路径，不接收完整数据数组；测试层为获得独立参考答案，可以持有完整输入及\code{std::sort}结果。二者的职责必须分清。

\section{课程要求及实现范围}
本报告以用户提供的《外存数据管理》课件、两份课堂语音转写和最终任务说明为课程依据。课件第75--81页给出Week 3--8推进顺序，分别覆盖文件与缓冲、归并计划、置换选择、动态缓冲、多路归并及测试评估。课堂说明要求最终报告附带有注释的源码，并能够解释重要函数和数据结构。

三项项目在同一仓库中逐步演进。矩阵文件示例用于理解Header、Directory和Data的组织思想；本实现将对象换为一维归并段。矩阵乘法、快排外排序、多块缓存替换和脏块机制不属于本次实现范围。这样的取舍使每一层都有明确输入输出，课堂演示可以直接追踪整数在文件、块和候选树之间的移动。

\section{Project 1：基于计划树的二路外排序}
Project 1首先解决可运行的完整外排序流程。输入是无文件头的连续本机字节序整数；输出是带目录的归并文件，以及说明最终段位置的\code{finalRun}。生成阶段每次读取不超过$M$个值，在内存中排序后按$B$写成一个有序段。由此得到至多$\lceil N/M\rceil$个初始段。

归并前先建立二叉计划。叶子引用已经存在的初始段，内部节点记录一次未来归并及两个孩子。计划构造与执行分开：先能打印整棵树、核对长度，再按后序递归执行。两路归并只保留两个输入块和一个输出块，不把两个完整段同时装入内存。奇数个节点时，最后一项带入下一层，避免无意义复制。

关键验收目标是：文件偏移正确、每个初始段及中间段有序、递归依赖正确、尾块不丢失、最终内容与参考排序逐项相同。仅显示“输出有序”不足以排除丢失或重复元素。

\section{Project 2：长段生成与置换选择}
普通分块生成的段长最多为$M$。若候选集合持续接纳仍能排在当前输出后面的新数，就有机会生成更长段，降低后续归并工作。Project 2维护固定$M$个候选，用败者树选出当前最小值，每输出一个再补入一个。若补入值小于刚输出值，则冻结，留给下一段；当前有效候选全部冻结时换段并解冻。

长段长度不等，因此归并顺序有实际影响。采用每次合并最短两段的哈夫曼计划，优化加权元素搬移量。并行发生于置换选择：输入线程预读备用块，计算线程操作活动块和树，输出线程写上一阶段的输出块。每阶段用\code{join}等待全部任务完成，再交接角色；后续二路归并仍为串行。

本项目必须体现真实败者数组、冠军路径更新、冻结状态、双输入/双输出以及段边界处理。不能以对全体数据排序替代置换选择，也不能仅启动空线程来声称完成并行。

\section{Project 3：K路归并与动态缓冲调度}
Project 3复用置换选择的初始段，归并阶段同时比较多路头值，使每轮段数按约$k$倍缩减。输入块不能永久固定为每路两个：各路消耗速度取决于关键字大小，某一路可能持续获胜，另一条路暂时不消耗。共享$2k$个物理块、各路队列和空闲池使空闲容量能够转给更需要的路。

预取依据是每路“最后一个已读关键字”，并在仍有未读数据的路中选择其最小者；缺少头值的空队列先补齐。未知头值意味着无法证明其他路当前值最小，因此必须在同步点等待补块。输出仍是一路，使用两块交替。多轮主控要正确处理不足$k$的末组、单段直接带过、空输入以及最终输出排空。

\begin{table}[htbp]\centering\caption{功能需求与验收依据}\label{tab:requirements}
\small\begin{tabularx}{\linewidth}{lYY}\toprule
需求 & 实现约束 & 可核验结果\\\midrule
文件组织 & 固定头、预留目录、流式追加 & 重开后偏移及数量正确，Scan/Put/Search有效\\
块缓冲 & 逐元素接口、尾块显式刷新 & 10个值、$B=3$时各4次数据块读写\\
二路计划 & 先建树、后序执行 & 15值示例4段、3次归并、$C=30$\\
置换选择 & $M$候选、冻结、败者路径重赛 & 15值示例段长7/7/1\\
最优计划 & 每次合并最短两段 & 7/7/1的代价由29减为23\\
阶段线程 & 三任务独占数据、join交接 & 输入/计算/输出有真实任务，最后一块排空\\
动态K路 & $2k$共享池、每路队列、跨路回收 & 测试中一路可占4块，回收块供另一条路使用\\
正确性 & 完整内容与参考排序一致 & 220次精确比较及机制断言通过\\
可演示性 & 原生构建、具名启动项、断点 & 一键十项验收与VS Code实测\\\bottomrule
\end{tabularx}\end{table}
